% Chapter 8 — Conclusion.
\section{Conclusion}

Hallucination detection needs to be cheap, fast, and explainable before it is useful in practice. This paper presented HaluRISC, a lightweight black-box risk predictor that combines 26 engineered evidence-consistency features with a tuned XGBoost classifier, calibrated probabilities, and SHAP explanations. On HaluEval QA it reaches an F1 of 0.9846 and an AUROC of 0.9979 at an expected calibration error of 0.0045, runs in 61.8\,ms per analysis at the median, and costs about \$0.001 per 1{,}000 predictions, roughly 100 times less than a GPT 5.6 Luna judge that it outperforms on the same 200 samples. A target-domain recalibration on 5{,}034 RAGTruth rows cuts the ECE on that corpus from 0.8185 to 0.1335, which shows that the calibration stage pays off exactly where the source-domain result said it would not.

The negative results are reported with the same care as the positive ones. The raw model needs no recalibration on HaluEval, the entity and hedging features add no measurable value, and zero-shot transfer to natural RAGTruth responses fails, with an F1 of 0.3022 on the QA split. The paper therefore frames the supervised risk model as a source-domain estimator, and the deployed system adds claim-level verification with evidence retrieval to cover the transfer gap. Future work targets multicalibration under distribution shift, cross-domain evaluation on FaithBench \cite{faithbench}, and neural baselines for comparison.
